\section{Design}
\label{sec:design}

This section presents the \pact type system and the gradual validation mechanism
that unifies type generation, workflow generation, and policy generation. We
describe the abstract syntax, the correspondence between \pact types and
JavaScript types, and the three forms of gradual validation that check
LLM-produced types, programs, and authorization policies at runtime.
\sect{sec:semantics} gives the operational semantics and soundness theorem.

\subsection{Syntax}
\label{sec:design-syntax}

\begin{sloppypar}
Figure~\ref{fig:syntax} presents the abstract syntax. A program is pragmas
followed by declarations. Types include base types (now including
\texttt{Principal} and \texttt{Path}), \texttt{Schema}, arrow types,
\texttt{Policy}, and named references. Expressions add to standard imperative
constructs the four primitives, \texttt{eval}, \texttt{enforce}, comparisons,
and the principal introducers \texttt{root} and \texttt{restrict}\;$e$. The
\texttt{as} clause on \texttt{gen}/\allowbreak\texttt{agent} takes an expression
of type \texttt{Principal}; the legacy literal-string form is dropped in favor
of named principals (\texttt{let user: Principal = restrict root}).
\end{sloppypar}

\begin{figure}[t]
\[\begin{array}{rcll}
P & ::= & \pi^*\; d^* & \textit{program} \\[4pt]
\pi & ::= & \texttt{@policy}\; s & \textit{policy pragma} \\
    & \mid & \texttt{@retries}\; n & \textit{retry pragma} \\[4pt]
d & ::= & \texttt{workflow}\; f(\overline{x_i{:}T_i}) \to T\; \{e\} & \textit{workflow declaration} \\
  & \mid & \texttt{type}\; X = T & \textit{type declaration} \\[4pt]
T & ::= & () \mid \texttt{String} \mid \texttt{Number} \mid \texttt{Bool} & \textit{base types} \\
  & \mid & \texttt{Schema} & \textit{type of types} \\
  & \mid & \texttt{Policy} & \textit{authorization policy} \\
  & \mid & (T_1 \to T_2) & \textit{arrow type} \\
  & \mid & \texttt{array[}T\texttt{]} & \textit{array type} \\
  & \mid & \{ \overline{x_i{:}T_i} \} & \textit{record type} \\
  & \mid & X & \textit{named type reference} \\[4pt]
e & ::= & s \mid n \mid b & \textit{literals} \\
  & \mid & x & \textit{variable reference} \\
  & \mid & \texttt{let}\; x {:} T = e \mid \texttt{let}\; x = e & \textit{immutable binding} \\
  & \mid & \texttt{var}\; x {:} T = e \mid \texttt{var}\; x = e & \textit{mutable binding} \\
  & \mid & x = e & \textit{assignment} \\
  & \mid & \texttt{if}\; (e)\; \{e\}\; [\texttt{else}\; \{e\}] & \textit{conditional} \\
  & \mid & \texttt{while}\; (e)\; \{e\} & \textit{loop} \\
  & \mid & \texttt{for}\; x\; \texttt{in}\; e\; \{e\} & \textit{iteration} \\
  & \mid & \texttt{ask}\; s & \textit{human prompt} \\
  & \mid & \texttt{say}\; s & \textit{human output} \\
  & \mid & \texttt{gen}\; s\; [\texttt{as}\; (s \mid x)] & \textit{LLM gen call} \\
  & \mid & \texttt{agent}\; s\; [\texttt{as}\; (s \mid x)] & \textit{LLM agent call} \\
  & \mid & \texttt{js}\; c & \textit{JavaScript expression} \\
  & \mid & \texttt{eval}\; e\;(e_1, \ldots, e_k) & \textit{workflow invocation} \\
  & \mid & \texttt{enforce}\; x & \textit{policy activation} \\
  & \mid & e_1 \;\mathit{op}\; e_2 & \textit{comparison} \\
  & \mid & e_1 ;\; e_2 & \textit{sequencing} \\[4pt]
\mathit{op} & ::= & \texttt{==} \mid \texttt{!=} \mid \texttt{<} \mid \texttt{<=} \mid \texttt{>} \mid \texttt{>=} & \textit{comparison operators}
\end{array}\]
\caption{Abstract syntax of Pact. Optional components are in square brackets and $\overline{x_i{:}T_i}$ denotes a possibly empty sequence.}
\label{fig:syntax}
\end{figure}

\begin{sloppypar}
\pact's base types correspond to JavaScript types
($\texttt{String}{\to}\texttt{string}$, etc.), and all runtime values are
represented as JSON. The type checker uses inference with subsumption.
\texttt{gen} is oracular, requiring explicit annotations at binding sites. Three
types play a special role in gradual validation: \texttt{Schema} is the
\emph{type of types}, arrow types are the types of workflows, and
\texttt{Policy} is the type of authorization rules. All three share the property
that their concrete values are generated by the LLM at runtime.
\end{sloppypar}

\subsection{Gradual Validation}
\label{sec:design-gradual}

The three gradual-validation annotations, namely \texttt{Schema},
arrow types, and \texttt{Policy}, share the same mechanism. Static
checking defers structural constraints on these types until the LLM
produces a concrete value at runtime, inspired by gradual
typing~\cite{siek2006gradual}. Rather than inserting
casts, \pact{} re-invokes the type checker over downstream code when a
concrete value materializes.

The common pattern is that an LLM-produced artifact is first accepted
only at an abstract type and is later forced when the program reaches a
use site that depends on its concrete structure. At that point, \pact{}
does not merely check the artifact in isolation. It re-checks the
program region that depends on the artifact. If the check fails, the
runtime records the violated constraint and retries the generator that
produced the artifact. This makes the retry causal. A missing field in a
generated schema retries the schema generator, not each later value
generator that happens to expose the missing field. An ill-typed
generated workflow retries the workflow generator, not the caller. A
malformed generated policy retries the policy generator before the
policy is installed. Our base assumption is that the LLM is \emph{eventually
truthful} in the sense that it will produce a correct artifact within a bounded
number of retries (\defn{def:oracle-truthful}).

\paragraph{Gradual type generation.} A \texttt{gen} call producing a
\texttt{Schema} type produces a JSON Schema at runtime~\cite{jsonschema2020}. 
When a schema materializes, the
interpreter converts it to the corresponding \pact{} type, registers it, and
re-invokes the type checker over the remainder of the program. On failure
(within \texttt{@retries}) the runtime feeds the type error back to the causal
\texttt{gen}. See Appendix~\ref{app:gradual-detail} for the formalization
of the forward syntactic analysis and validation procedures.

\paragraph{Gradual workflow generation.} A \texttt{gen} call producing an arrow
type instructs the LLM to produce a complete \pact{} workflow as structured
output. Syntactic discrepancies are build-in and we can assume a \pact workflow will
be well-formed, but not necessarily well-typed. Workflow gradual validation then
takes effect, substituting the generated workflow at the call site (\texttt{eval}) and re-checking the caller's
continuation. On failure, the runtime retries the workflow generator with the error as feedback,
same as with gradual type generation. This is a powerful mechanism for
enforcing compositionality and modularity of LLM-generated workflows, which are otherwise black boxes. 
For example, in Figure~\ref{fig:overview-workflow} the workflow generator produces a subroutine that is called three times. 
If the generated workflow is ill-typed, the repair happens once at the source,
and all three calls conform on their first attempt. 

\paragraph{Gradual policy generation.} A \texttt{gen} call producing a
\texttt{Policy} type instructs the LLM to produce a Cedar policy as JSON. \pact
allows first class \emph{principals}---defined from others using the \texttt{restrict} command---
and \emph{policies}, generated by \texttt{gen} and activated by an \texttt{enforce} command later.
Separating generation from
activation lets a program generate a candidate policy, inspect it,
branch on it, or route it through human approval before it affects later
calls. A failed policy compile or authorization check triggers
feedback-driven self-healing.
Policy gradual validation checks that the generated policy is well-formed and
references only declared principals after it is enforced.
A failed validation retries the policy generator with
the error as feedback, before the policy is installed.

\paragraph{Cedar policies in \pact{}.} Cedar requests in \pact{} are
\texttt{(principal, action, resource)} triples with $\texttt{Action}\in
\{\texttt{Read},\texttt{Write},\texttt{Exec}\}$ and \texttt{Resource}
ranging over files and directories the agent ultimately touches;
wildcards \texttt{*} are supported in either slot. Principals form a
tree rooted at the top-level principal \texttt{root}; the policy
state begins as $\MPol_0 = \texttt{allow(root,\;*,\;*)}$ root is allowed
indiscriminate access to all files in the system, much like in the POSIX
model. Forbid policies are also allowed, for example \texttt{forbid(user, "Write", "/")}
prohibits the principal \texttt{user} from writing to any file. The policy state
is updated by appending new rules with the \texttt{enforce} command, and the
decision procedure is monotone in the deny direction.
The only way to create new principals is via \texttt{restrict}, which
creates a fresh child of an existing principal. 
$\texttt{let}\;b = \texttt{restrict}\;a$ declares $b$ as a fresh
descendant of $a$, inheriting all of $a$'s allowances by Cedar's
hierarchy semantics, for example one can create
a first principal by doing $\texttt{let}\;coordinator = \texttt{restrict}\;\texttt{root}$,
as seen in Figure~\ref{fig:clinical-trial}. Then a new policy
\texttt{enforce} admits only candidate policies in
which (i) every rule is forbid-effect, (ii) only principals
declared at the time the policy is enforced are referenced, and (iii) no rule references \texttt{root}.
These rules are required to prove the \emph{policy monotonicity} property of the runtime semantics, which states that the allow-set can only shrink across steps 
(Theorem~\ref{thm:pact-soundness}).
Because no rule can reference \texttt{root} and no rule can be allow-effect, no
rule can expand the authority of any principal, and because new principals are
created as descendants of existing ones,
no rule targets
\texttt{root}; together these admit only further restriction of
descendants from their inherited authority while leaving \texttt{root}
permanently maximal.

The static premise enforced at every \texttt{gen}/\texttt{agent} redex
is exactly $\PrinDeclared{\MPol}{\MPrin}$: the principal named in the
\texttt{as} clause, or \texttt{root} when none is given, is declared in
the current policy state. The action and resource of any specific
filesystem access become known only when the agent runs, so they are
checked just-in-time by the runtime monitor of
\sect{sec:implementation}.

\subsection{Threat Model and Limitations}
\label{sec:design-threat}

\pact{} trusts the programmer-authored text, the type checker, the
runtime re-checker, and the Cedar decision procedure, and treats every
LLM output as untrusted. The retry budget $\MRetry$ bounds how many
untrusted attempts the runtime consumes per site. Validation rejects
outputs that fail structural checks of the annotated type. Every
\texttt{gen}/\texttt{agent} carries the static premise
$\PrinDeclared{\MPol}{\MPrin}$ before the LLM is called, and every
filesystem access the agent then performs is policy-checked
just-in-time by the runtime monitor (\sect{sec:implementation}).
A type-validation failure triggers retry; a policy denial fails closed.

\paragraph{Authority is monotone in the deny direction.} Because
\texttt{enforce} installs only forbid-effect rules on non-root
principals (\sect{sec:design-gradual}), the allow-set
$\{(\MPrin, a, r) \mid \mathsf{allows}(\MPol, \MPrin, a, r)\}$
monotonically shrinks across every step. Generated policies can
therefore only restrict authority of declared descendants, never
expand it; this rules out privilege escalation through LLM-authored
policies as a class.

\label{sec:design-limitations}
Four threat classes remain outside this model.
These are prompt injection, semantic misalignment of generated policies,
\texttt{say}-based exfiltration, and host-side JavaScript effects.
They are expanded in \sect{app:threat}. Formal results are stated
relative to an \emph{eventually-truthful} oracle
(\defn{def:oracle-truthful}). An adversarial or permanently lying oracle
exhausts the budget and the configuration fails closed.
\sect{app:limitations} enumerates each restriction.
